\scanpage{56}
% manuscript line 1
$h$ compact et la série \inserted{$(\rho_i)$} de ses valeurs propres convergente, on aura
$h=\sum\rho_i\bar a_i\otimes a_i$, où $(a_i)$ est orthonormal \deleted{\unsure{est}}, d'où
$\|h\|_1\leq\sum\rho_i\|a_i\|^2=\sum\rho_i$, ce qui prouve que $h$ est op. de
Fredholm, et qu'on a (20). On peut d'ailleurs remplacer
\deleted{l\unsure{e} p\unsure{r.}} $h=\sqrt{u^*u}$ par $h'=\sqrt{uu^*}$, \deleted{cel\unsure{a se vo}it} soit
directement, soit par application du fait \deleted{\unsure{suivant}} \inserted{bien connu} \deleted{que si $u$ est}
\deleted{\unsure{et $h$}} \uline{un op. compact quelconque} \inserted{$u\in L(E,F)$} \uline{alors} $\sqrt{u^*u}$ et $\sqrt{uu^*}$
\uline{ont les mêmes valeurs propres}. En effet, si $(\lambda_i)$ sont
les valeurs propres de $h=\sqrt{u^*u}$, on aura $h=\sum\rho_i\bar a_i\otimes a_i$
et $u=Uh$, où $U$ est isométrique sur l'espace engendré
par les $a_i$, d'où $U$$=\sum\rho_i\bar a_i\otimes b_i$, où $(b_i)=(Ua_i)$
forment un système orthonormal \inserted{dans $F$}. On en conclut donc

% manuscript line 13
\[
\begin{aligned}
u^*&=\sum\rho_i\bar b_i\otimes a_i,\quad\text{et}\quad uu^*=\sum_{ij}\rho_i\rho_j(a_i,a_j)\bar b_j\otimes b_i\\
&=\sum_i\rho_i^2\bar b_i\otimes b_i\,;\quad\text{d'où}\qquad
\sqrt{uu^*}=\sum_{\unsure{1|i}}\rho_i^{\unsure{2|1}}\bar b_i\otimes b_i,
\end{aligned}
\]

% manuscript line 15
\uline{Corollaire 1} Les op. de Fredholm de $E$ dans $F$ \deleted{\unsure{sont ceux de}} \inserted{sont \deleted{ceux} les}
opérateurs compacts de la forme

% manuscript line 16
\[
\boxed{u=\sum\rho_i\bar a_i\otimes b_i} \tag*{\textcircled{21}}
\]

% manuscript line 17
où $(a_i)$ (resp. $(b_i)$) est un système orthonormal
dans $E$ (resp. $F$), et $(\rho_i)\in l^1$. On a $\|u\|_1=\sum\rho_i$.

% manuscript line 19
% margin note written sideways
\margin{L'espace des applications de Fredholm de $E$ dans $F$ est noté $L^{(1)}(E,F)$, on a donc $L^{1}(E,F)=\bar E\widehat\otimes F=E'\widehat\otimes F$. \unsure{car} $E'=\bar E$.}

% manuscript line 19
Comme $\unsure{\bar E|E}'$$\widehat\otimes E\to L(E)$ est biunivoque, on peut parler
de la trace d'un op. de Fredholm dans l'espace de Hilbert $E$.
\deleted{\unsure{On aura donc}} La formule (20) peut s'écrire aussi

% manuscript line 22
\[
\boxed{\|u\|_1=\operatorname{Tr}\sqrt{u^*u}=\operatorname{Tr}\sqrt{uu^*}} \tag*{\textcircled{22}}
\]

% manuscript line 23
\uline{Corollaire 2} Si $u$ est un op. de Fredholm de $E$ dans $F$,
on a

% manuscript line 25
\[
\|u\|_1=\sup_{\|A\|\leq1}|\operatorname{Tr} A\circ u|\qquad(A\in L(F,E)) \tag*{\textcircled{22}}
\]

% manuscript line 26
(En effet, le dual de $E'$$\widehat\otimes F$ est $B(E'',F)\simeq L(F,E)$) \inserted{\unsure{\ldots}v. N. \unsure{ant.}}

% manuscript line 27
% margin note written sideways
\margin{Conjuguant avec le th. 2, formule (19), on trouve le \uline{Corollaire 3} Pour que $u$ soit un op. de H.-S., il faut et il suffit que $u^*u$ soit un op. de Fredholm, et alors on a $(23)\quad\|u\|_2=\sqrt{\operatorname{Tr}(u^*u)}\deleted{\leq}=\sqrt{\|u^*u\|_1}$.}

% manuscript line 28
\deleted{\unsure{On}} Les relations entre opérateurs de Fredholm et op.
de H.-S. sont incluses dans l'important

% manuscript line 30
\uline{Th. 4} Soient \deleted{$E$ \unsure{et} $F$\unsure{,} d\unsure{eux}} \inserted{$E,F,G$ des} espaces de Hilbert, $u\in L^2(E,F)$, $v\in L^2(F,G)$,
Alors $v\circ u\in L^1(E,G)$, et \deleted{o\unsure{n}}

% manuscript line 32
\[
\boxed{\|vu\|_1\leq\|v\|_2\,\|u\|_2} \tag*{\textcircled{24}}
\]

% manuscript line 33
% archive sheet number: 57
